\begin{center}\begin{tikzpicture}[font=\small,>=stealth,line join=round]
\draw[gray!40,rounded corners=3pt] (0,0) rectangle (14.2,8.3);\node at (7.1,7.700000000000001) {例8：同一横坐标既有轴又有中心，函数只能恒为4};
\begin{scope}[shift={(2.3857142857142857,2.003846153846154)},x=0.8571428571428571cm,y=0.7076923076923076cm]\draw[->,gray] (-1.5,0)--(5.5,0) node[right]{$x$};\draw[->,gray] (0,-0.5)--(0,6) node[above]{$y$};\draw[blue,thick,solid,domain=-1:5,samples=180] plot(\x,{4});\fill[orange!85!black] (2,4) circle(2.2pt);\draw[violet,dashed] (2,-.3)--(2,5.6);\end{scope}\node[anchor=west,align=left,text width=5.0cm] at (8.4,4.0) {轴 $x=2$\\[9pt]中心 $(2,4)$\\[9pt]两值相等，又相加为8\\[9pt]每个函数值都是4\\[9pt]$f\equiv4$};
\node[text width=12.6cm,align=center] at (7.1,.65) {同一竖直线上的两种要求同时成立，会使整条图像变为水平线。};
\end{tikzpicture}\end{center}
